\section{Related Work}

\subsection{Water Quality Prediction}
Data-driven prediction of water-quality indicators has been widely explored with ANN, tree-based models, and recurrent architectures. Reviews and recent studies report consistent use of physicochemical variables and geospatial context for forecasting pH, OD, turbidity, and composite indicators \cite{yan2024,ahmed2019,aladhaileh2021,pan2025}. LSTM-based approaches are common where temporal records are dense \cite{liu2019,pyo2023,qi2020}, while hybrid and ensemble methods improve robustness under nonlinear behavior and limited sample regimes \cite{wu2022,chadalavada2025,kalaivanan2022,palani2008}.

Beyond model choice, an important distinction in the literature is the operational objective: some studies prioritize maximum predictive accuracy in centralized infrastructure, while others prioritize robustness under noisy sensing, sparse sampling, and field constraints. For inland-water applications, this distinction is critical because model behavior is affected by local heterogeneity, seasonality, and variable measurement uncertainty \cite{yan2024,lowe2022,rahu2023}.

\subsection{Edge AI and Energy Evaluation}
For environmental IoT deployments, inference cost must be evaluated jointly with predictive quality. Prior work highlights the relevance of platform-level benchmarking and hardware-aware model design \cite{rosero2020,zulkifli2022,lamrini2024,elshafeiy2023}. Recent neuromorphic and energy-efficient AI studies indicate that event-driven models can be attractive for constrained settings, but deployment evidence across heterogeneous edge hardware is still developing \cite{aral2024,azadi2025,huijbregts2024,aribe2025future}. This gap motivates the present work: a unified comparison of offline multimodal inference quality and energy-per-inference behavior on multiple practical platforms.

Most prior studies report either algorithmic metrics or hardware metrics in isolation. In contrast, field deployment requires a coupled view of latency, throughput, and incremental energy under realistic idle/load transitions. Consequently, cross-platform benchmarking should be interpreted as a systems result, not only as a model result.
